\section{The flat amplitude prior, written in terms of the reference frequency}
\label{sec:flat}

That the flat-prior marginal of Eq.~(\ref{eq:Lflat}) pulls spectral indices is known.
\citet{eriksen2008} found a global offset biased by more than $4\sigma$ under uniform priors and
removed it with a Jeffreys prior on the index. \citet{dunkley2009}, fitting one polarized synchrotron
index per $N_{\rm side}=2$ pixel to the five-year WMAP maps with uniform priors on the CMB and
synchrotron amplitudes defined at 23 GHz, found index posteriors centred on $-2.6$ in simulations,
noted that the effect grows with the number of pixels that share an index and that a Jeffreys prior
on the index does not remove it, and divided their prior by a factor estimated from a noise-only
simulation. \citet{alonso2017} cancel the determinant with a ``volume prior'', which in our reading
makes their target the profile likelihood, and show biases of more than $3\sigma$ without it. Jeffreys
or reference priors are used against prior-volume effects in several other analyses
\citep{dickinson2009,planck2015x,jew2019}, and \citet{andersen2023} describe the marginal and the
profile (``ridge'') index samplers of Commander, noting that the two differ by the determinant term. What has not been
written down, as far as we can find, is how large the pull is in closed form and what sets its sign.
Both follow from one identity.

\subsection{An exact identity}

Change the reference frequency of the power-law component from $\nuref$ to $\nuref'$. The mixing
column of every amplitude in index region $r$ is multiplied by $(\nuref/\nuref')^{\beta_r}$, so
$\M' = \M\,\mathsf{D}$ with $\mathsf{D}$ diagonal, $\chimin$ is unchanged, and
$\det\A' = \det\A \prod_r (\nuref/\nuref')^{2 n_r \beta_r}$, with $n_r$ the number of power-law
amplitudes in region $r$. Therefore
\begin{equation}
  \ln \mathcal{L}_{\rm flat}(\beta;\nuref') = \ln \mathcal{L}_{\rm flat}(\beta;\nuref)
  + \sum_r n_r\, \beta_r \ln\frac{\nuref'}{\nuref} ,
  \label{eq:tilt}
\end{equation}
for any number of components, any bandpasses and any noise covariance. Moving the frequency at which
a flat amplitude prior is stated is the same as multiplying the index prior by an exponential,
$e^{\lambda\beta}$ with $\lambda = n \ln(\nuref'/\nuref)$: 34 $e$-folds per unit of index for 128
amplitudes moved from 23 to 30 GHz. The profile likelihood and the zero-mean prior with a free scale
do not change.

\subsection{The shift and its pivot}

Let $g_r = -\tfrac12\,\partial \ln\det\A/\partial\beta_r$ be the gradient of the volume term. To first
order the maximum of $\mathcal{L}_{\rm flat}$ sits at the profile maximum plus
\begin{equation}
  \delta\beta = \mathsf{H}^{-1} g , \qquad
  \mathsf{H} = \J + \nabla\nabla\,\tfrac12\ln\det\A ,
  \label{eq:response}
\end{equation}
where $\mathsf{H}$ is the expected curvature of the flat-prior marginal itself: the profile information
plus the curvature of the volume term. For one power-law component with delta bands and diagonal
noise, $-\tfrac12 \ln A_p = \beta\ln\nuref - \tfrac12 \ln\sum_b w_{pb}\, \nu_b^{2\beta}$, so
\begin{equation}
  g = \sum_p \ln\frac{\nuref}{\bar\nu_p}, \qquad
  \ln\bar\nu_p = \frac{\sum_b w_{pb}\,\nu_b^{2\beta} \ln\nu_b}{\sum_b w_{pb}\,\nu_b^{2\beta}} , \qquad
  \nabla\nabla\,\tfrac12\ln\det\A = 2\sum_p V_p = 2\,\Tm ,
\end{equation}
where $\bar\nu_p$ is the frequency at which amplitude $p$ is best measured: the geometric mean of the
band frequencies under the weights (inverse noise variance times squared spectrum) that also define
the lever arm $V_p$ in Eq.~(\ref{eq:uV}). When the amplitudes of a region share a noise pattern,
$\bar\nu_p = \nupiv$ and, with $\sigma_\beta^2 = 1/\J$ the profile variance of the region index,
\begin{equation}
  \boxed{\;\delta\beta = \ln\!\left(\frac{\nuref}{\nupiv}\right) \frac{n\,\sigma_\beta^2}{1 + 2\,\Tm/\J}
         \;\longrightarrow\; \ln\!\left(\frac{\nuref}{\nupiv}\right) \sigma_1^2
         \quad (\Tm \ll \J)\;}
  \label{eq:shift}
\end{equation}
where $\sigma_1^2 = n\sigma_\beta^2$ is the index variance that a single amplitude carries, and
$\Tm/\J$ is the inverse mean squared signal-to-noise ratio of the amplitudes (Sect.~\ref{sec:profile}).
The boxed form is for one component. With further free amplitudes the curvature of the volume term
is no longer $2\Tm$ and can be negative, so that $\mathsf{H}$ is smaller than $\J$ and the flat-prior
marginal may have no interior optimum at all; one of our fits to real maps, for a band set without K band
and with a free CMB amplitude (on the earlier, unmasked regions described in Sect.~\ref{sec:data}), ends
at the edge of the allowed range (\exref{E040}). Three things follow.

\emph{The shift does not average down.} Adding amplitudes at fixed depth leaves $\sigma_1$ and
$\Tm/\J$ unchanged, so $\delta\beta$ stays the same while the posterior width falls as $n^{-1/2}$. In
units of the posterior width the pull grows as $\sqrt{n}$: this is the observation of
\citet{dunkley2009} that the effect increases with the number of pixels per index, and it is the
Neyman-Scott structure \citep{neyman1948}. More data of the same kind make it worse. So does a finer
pixelization of the same data: for $\Tm \ll \J$, four times as many amplitudes for the same region
information is four times the shift.

\emph{The sign is set by the convention.} The index is pulled steep if the reference lies below the
pivot and flat if above. A flat prior is neutral only when stated at the pivot.

\emph{The pivot depends on the other components and on their priors.} A component with a free
amplitude and a spectrum that does not depend on the index (the CMB, dust with fixed parameters)
contributes nothing to $g$ directly, but projecting it out replaces the weights $w_{pb}$ by the
projected weight matrix $\tilde{\mathsf{W}}_p = \mathsf{W}_p - \mathsf{W}_p \mathsf{F}
(\mathsf{F}\T\mathsf{W}_p\mathsf{F} + \mathsf{S}^{-1})^{-1}\mathsf{F}\T\mathsf{W}_p$, with $\mathsf{F}$ the spectra of the
other components and $\mathsf{S}$ their prior covariance. The pivot becomes a generalized mean under
$\tilde{\mathsf{W}}_p$, which has negative off-diagonal elements and can lie outside the range of the
bands. With a free CMB amplitude under a flat prior the index information sits in the lowest band and
the pivot moves to or below it. With a CMB prior the size of the signal it moves back.

For several regions, several components and dense noise, $g$ and $\mathsf{H}$ are a vector and a
matrix, and Eq.~(\ref{eq:tilt}) gives the difference between two references,
\begin{equation}
  \hat\beta_{\rm flat}(\nuref') - \hat\beta_{\rm flat}(\nuref) \simeq
  \mathsf{H}^{-1}\,\big[\,n_r \ln(\nuref'/\nuref)\,\big] ,
  \label{eq:coupled}
\end{equation}
in which the regions are coupled: a large region with many amplitudes drags its neighbours through
the off-diagonal elements of $\mathsf{H}^{-1}$. We test this form on simulations and on real data in
Sect.~\ref{sec:data}. Equation~(\ref{eq:response}) was corrected by an independent review of our
first derivation, which had $\J$ in place of $\mathsf{H}$; the two agree where $\Tm \ll \J$, and a
direct test with the reference placed close to the pivot gives a measured shift of $1.01$ and $1.03$
times Eq.~(\ref{eq:response}) at $\Tm/\J = 0.03$ and $0.11$ where the uncorrected form is off by 5 and
16 percent (\exref{E039}, independent review). For shifts that are not small against the range over
which the curvature changes, neither is accurate (Fig.~\ref{fig:ladder}, right).

\subsection{Sizes at the depth of the public polarization data}

\begin{table}
\centering
\caption{The flat-prior marginal against Eq.~(\ref{eq:shift}) on simulations at the white-noise depth
of the WMAP and LFI polarization maps (\exref{E036}). One index per $N_{\rm side}=2$ pixel (128 $Q/U$
amplitudes at $N_{\rm side}=16$), exact one-dimensional posteriors on a grid, 100 noise realizations;
truth: the Cosmoglobe DR1 polarized synchrotron map smoothed to $5^\circ$ with $\beta = -3.1$. Regions
are grouped by their profile index error $\sigma_\beta$. ``Shift'' is the mean of posterior mean minus
truth. The predicted value is the high signal-to-noise limit of Eq.~(\ref{eq:shift}) with the class median of
$\sigma_\beta$ (in brackets) and the pivot of the band set; it is first order in the shift and uses one number for a class, so it
degrades where the shift exceeds the posterior width.}
\label{tab:flat}
\small
\begin{tabular}{llrrrr}
\toprule
model & $\sigma_\beta$ class (median) & reference & shift & predicted & shift / posterior sd \\
\midrule
synchrotron only, & below 0.03 (median 0.019) & 23 GHz & $-0.005$ & $-0.006$ & $-0.26$ \\
bands K, 30, Ka, Q, 44 &             & 30 GHz & $+0.007$ & $+0.006$ & $+0.35$ \\
($\nupiv = 26.4$ GHz)   &             & 60 GHz & $+0.037$ & $+0.038$ & $+1.9$ \\
                        & 0.03 to 0.1 (0.059) & 23 GHz & $-0.061$ & $-0.062$ & $-0.94$ \\
                        &             & 30 GHz & $+0.062$ & $+0.057$ & $+1.01$ \\
                        &             & 60 GHz & $+0.341$ & $+0.37$ & $+6.1$ \\
                        & 0.1 to 0.25 (0.161) & 23 GHz & $-0.40$ & $-0.46$ & $-2.4$ \\
                        &             & 30 GHz & $+0.34$ & $+0.42$ & $+2.5$ \\
\midrule
free CMB and dust, flat priors, & 0.03 to 0.1 & 23 GHz & $+0.006$ & & \\
adding V, 70, 353 GHz           &             & 30 GHz & $+0.167$ & & \\
                                & 0.1 to 0.25 & 23 GHz & $+0.058$ & & \\
                                &             & 30 GHz & $+0.859$ & & \\
\midrule
the same with a CMB prior & 0.03 to 0.1 & 23 GHz & $-0.060$ & & \\
of 0.2 $\muK$ per element &             & 30 GHz & $+0.063$ & & \\
\bottomrule
\end{tabular}
\end{table}

Table~\ref{tab:flat} gives the sizes on a realistic polarized sky at the white-noise depth of the
WMAP and LFI maps. For the synchrotron-only model the pivot of the five lowest bands is 26.4 GHz, and
references of 23 and 30 GHz, the two in use, sit on opposite sides of it: the same data give an
index steeper than the truth by 0.06 in one convention and flatter by 0.06 in the other, one
posterior standard deviation each way, for regions measured to 0.06. For regions measured to 0.16
the pull is 0.4 and 0.3. The difference between the two references follows Eq.~(\ref{eq:tilt}):
$\ln(30/23)\times128\times0.059^2 = 0.118$ against a measured 0.123. Doubling and quadrupling the
noise leaves the shift at a given $\sigma_\beta$ unchanged, as Eq.~(\ref{eq:shift}) requires
($-0.063$ and $-0.073$ at 23 GHz for the class 0.03 to 0.1; \exref{E036}). The profile likelihood on
the same simulations is unbiased to 0.002 in the first class and 0.02 in the second
(Sect.~\ref{sec:profile}).

With free CMB and dust amplitudes under flat priors the pivot moves to 22 to 23 GHz. A flat
prior stated at 23 GHz is then nearly neutral and one stated at 30 GHz pulls the index flat by 0.17
to 0.86. With a Gaussian CMB prior of 0.2 $\muK$ per element, the size of the CMB polarization signal
at this resolution, the pivot and the shifts return to the synchrotron-only values, and so does the
index error, which the flat CMB prior had inflated by a factor 1.63. The prior on the CMB amplitude
matters more for the synchrotron index than the prior on the synchrotron amplitude does, as long as
the latter is not flat.

\paragraph{An existing real-data comparison.} An earlier experiment of ours (\exref{E022}) evaluated
both the flat-prior marginal and the profile likelihood on real WMAP Ka to V, LFI and 353 GHz
polarization maps, in four regions reconstructed from the released index map of \citet{svalheim2023} and
without a mask (Sect.~\ref{sec:data} explains why we no longer use those regions; the argument here is
differential and does not depend on what they are), at a 30 GHz reference, for two noise models. In its high signal-to-noise limit Eq.~(\ref{eq:shift}) says the difference of the two estimates divided by
$n\sigma_\beta^2$ is $\ln(30/\nupiv)$ whatever the noise model. In the region reconstructed as the Spur the two noise
models give shifts of $+0.271$ and $+0.172$ for $n\sigma_\beta^2 = 1.913$ and $1.216$: a pivot of 26.04
GHz in both cases, to 0.2 percent, for a band set whose lowest band is at 28.4 GHz and with a free CMB
amplitude. With a CMB prior of 0.19 $\muK$ the implied pivot is 29.5 to 30.2 GHz. For that band set
the neutral reference frequency is therefore 26 GHz or 30 GHz depending on how the CMB is
constrained. For WMAP K to V on the real maps it is \N{pivot.range} GHz with the $\Lambda$CDM prior on
the CMB and about \N{scan.flat.pivot} GHz at high latitude, below the lowest band, with a flat one,
where the flat-prior marginal then has no optimum at all (Sect.~\ref{sec:data}).
Sect.~\ref{sec:data} tests the identity directly on the WMAP
maps with their dense covariances.
